\Subhead{3.6.2}{Флуктуационно-диссипационная теорема}

Одна и та же тепловая среда тормозит движение и создаёт флуктуации. В равновесии их интенсивности связаны так, чтобы сохранялось распределение Гиббса. Для осциллятора с силой, сопряжённой координате, $x(\omega)=\chi(\omega)F(\omega)$, $\chi=(k-m\omega^2-\ii b\omega)^{-1}$. Средняя поглощаемая мощность гармонической силы $F_0$ равна $\omega\chi''|F_0|^2/2$; поэтому $\chi''$ измеряет диссипацию.

Классическая ФДТ в нашей двусторонней нормировке:
\begin{equation}\label{eq:fdt}
 S_x(\omega)=\frac{2k_BT}{\omega}\chi''(\omega).
\end{equation}
Для осциллятора $\chi''=b\omega|\chi|^2$, откуда $S_F=2bk_BT$. Нормировка проверяется интегралом $\langle x^2\rangle=k_BT/k$; независимо задавать равновесное трение и шум нельзя.

Квантовая версия для \emph{симметризованной} корреляции:
\begin{equation}
 S_x^{\rm sym}(\omega)=\hbar\coth\frac{\hbar\omega}{2k_BT}\,\chi''(\omega).
\end{equation}
Она переходит в классическую при $\hbar\omega\ll k_BT$. Несимметризованные спектры положительной и отрицательной частоты различаются и описывают способность среды поглощать и отдавать кванты. Нулевые колебания в симметризованном спектре при $T=0$ не являются источником извлекаемой тепловой мощности. ФДТ требует теплового равновесия и линейного отклика; активный накачиваемый генератор обычно этим условиям не удовлетворяет.
